% GENERATED by scripts/gen_blueprint.py — do not edit.
\chapter{Directed paths: the Cheng--Keevash results}
\label{bp:ch_paths}

\section{The $\delta = 3$ path theorem}

\begin{theorem}[\code{ck\_conj1\_delta3}]
  \label{res:ck_conj1_delta3}
  \uses{def:arc, def:dipath, def:ell, def:oriented, def:outdeg}
  \rocq{Digraph.applications.ck3.ck3_main.ck_conj1_delta3}
  \rocqok

\begin{verbatim}
Theorem ck_conj1_delta3 (D : orientedDigraph) :
  0 < #|D| -> (forall v : D, 3 <= outdeg v) -> 6 <= ell D.
\end{verbatim}
\href{https://github.com/LLM4Rocq/digraph-theory/blob/main/theories/applications/ck3/ck3_main.v\#L147}{Source: \code{theories/applications/ck3/ck3\_main.v:147}}


\par\textbf{Decoded.} Let $D$ be any nonempty finite oriented digraph in which every vertex
has out-degree at least $3$. Then $\ell(D) \ge 6$
(\ref{def:ell}): $D$ contains a directed simple path with $6$ arcs
(7 distinct vertices). The second form unfolds $\ell$ into the
explicit existence of such a path (\ref{def:dipath}).
This is the $k = 3$ case of Cheng--Keevash Conjecture~1.

\end{theorem}

\begin{theorem}[\code{ck\_conj1\_delta3\_path}]
  \label{res:ck_conj1_delta3_path}
  \uses{def:arc, def:dipath, def:oriented, def:outdeg}
  \rocq{Digraph.applications.ck3.ck3_main.ck_conj1_delta3_path}
  \rocqok

\begin{verbatim}
Corollary ck_conj1_delta3_path (D : orientedDigraph) :
  0 < #|D| -> (forall v : D, 3 <= outdeg v) ->
  exists x : D, exists s : seq D, dipath x s /\ size s = 6.
\end{verbatim}
\href{https://github.com/LLM4Rocq/digraph-theory/blob/main/theories/applications/ck3/ck3_main.v\#L158}{Source: \code{theories/applications/ck3/ck3\_main.v:158}}


\end{theorem}

\section{The $\delta = 2$ case}

\begin{theorem}[\code{ck\_conj1\_delta2}]
  \label{res:ck_conj1_delta2}
  \uses{def:arc, def:dipath, def:ell, def:oriented, def:outdeg}
  \rocq{Digraph.applications.ck3.lemma7.ck_conj1_delta2}
  \rocqok

\begin{verbatim}
Corollary ck_conj1_delta2 (D : orientedDigraph) :
  0 < #|D| -> (forall v : D, 2 <= outdeg v) -> 4 <= ell D.
\end{verbatim}
\href{https://github.com/LLM4Rocq/digraph-theory/blob/main/theories/applications/ck3/lemma7.v\#L698}{Source: \code{theories/applications/ck3/lemma7.v:698}}


\par\textbf{Decoded.} Same statement one level down: minimum out-degree $\ge 2$ forces a
directed simple path of length $4$.

\end{theorem}

\section{Theorem 4, oriented version}

\begin{theorem}[\code{ck\_theorem4\_oriented}]
  \label{res:ck_theorem4_oriented}
  \uses{def:arc, def:dipath, def:ell, def:oriented, def:outdeg}
  \rocq{Digraph.applications.ck3.lemma7.ck_theorem4_oriented}
  \rocqok

\begin{verbatim}
Theorem ck_theorem4_oriented (D : orientedDigraph) (delta : nat) :
  0 < #|D| -> (forall v : D, delta <= outdeg v) ->
  2 * delta - (delta - 1)./2 <= ell D.
\end{verbatim}
\href{https://github.com/LLM4Rocq/digraph-theory/blob/main/theories/applications/ck3/lemma7.v\#L684}{Source: \code{theories/applications/ck3/lemma7.v:684}}


\par\textbf{Decoded.} Every nonempty oriented digraph with minimum out-degree $\ge \delta$
satisfies $\ell(D) \ge 2\delta - \lfloor(\delta-1)/2\rfloor$
(\code{./2} is truncated halving). For even $\delta$ this is
sharper than the $\lceil 3\delta/2 \rceil$ stated by Cheng--Keevash.

\end{theorem}

\section{Cheng--Keevash Lemma 7, uniform in $\delta$}

\begin{theorem}[\code{lemma7}]
  \label{res:lemma7}
  \uses{def:arc, def:dipath, def:ell, def:oriented, def:outdeg}
  \rocq{Digraph.applications.ck3.lemma7.lemma7}
  \rocqok

\begin{verbatim}
Theorem lemma7 (D : orientedDigraph) (delta : nat) :
  0 < #|D| -> (forall v : D, delta <= outdeg v) ->
  2 * delta <= ell D \/
  exists S : {set D},
    [/\ S != set0, #|S| <= delta
      & forall v, v \in S -> 2 * delta - ell D <= outdeg_in S v].
\end{verbatim}
\href{https://github.com/LLM4Rocq/digraph-theory/blob/main/theories/applications/ck3/lemma7.v\#L646}{Source: \code{theories/applications/ck3/lemma7.v:646}}


\par\textbf{Decoded.} In a strong (\ref{def:strong}) oriented digraph that is
$\delta$-out-regular-from-below ($d^+(v) \ge \delta$ everywhere,
$\delta \ge 1$) and where the longest path is short
($\ell(D) \le 2\delta$), the vertex set of any longest path supports
the structural ``kernel'' conclusion quantified in the statement ---
the engine behind the path theorems. (This page quotes the statement
in full; the surrounding section fixes the oriented digraph
\code{D}, $\delta$, and the strongness/degree hypotheses
visible in the quoted text.)

\end{theorem}

\section{No short-path strong counterexample at $\delta = 3$}

\begin{theorem}[\code{no\_short\_strong3}]
  \label{res:no_short_strong3}
  \uses{def:arc, def:dipath, def:ell, def:oriented, def:outdeg, def:strong}
  \rocq{Digraph.applications.ck3.ck3_main.no_short_strong3}
  \rocqok

\begin{verbatim}
Lemma no_short_strong3 (H : orientedDigraph) :
  (forall v : H, outdeg v = 3) -> 0 < #|H| -> strongb H ->
  ell H <= 5 -> False.
\end{verbatim}
\href{https://github.com/LLM4Rocq/digraph-theory/blob/main/theories/applications/ck3/ck3_main.v\#L28}{Source: \code{theories/applications/ck3/ck3\_main.v:28}}


\par\textbf{Decoded.} There is no strong oriented digraph with $\delta^+ \ge 3$ and
$\ell \le 5$: the reduction target of the main theorem, stated
separately.

\end{theorem}

